\section{Week 7: Model Adaptation, LoRA, Quantization, and Deployment Reality}

\subsection{Learning Objectives}

By the end of this week, you should be able to:
\begin{itemize}
  \item Explain why full fine-tuning is rarely used in production
  \item Understand parameter-efficient fine-tuning (PEFT)
  \item Describe LoRA and how it modifies linear layers
  \item Reason about quantization and precision tradeoffs
  \item Design a validation strategy for adapted models
\end{itemize}

\subsection{LLM Components Covered}

\textbf{Adaptation}
\begin{itemize}
  \item Fine-tuning vs parameter-efficient fine-tuning
  \item Low-Rank Adaptation (LoRA)
\end{itemize}

\textbf{Deployment}
\begin{itemize}
  \item Quantization
  \item Memory and latency constraints
  \item Model versioning
\end{itemize}

\subsection{Conceptual Core: Why Full Fine-Tuning Breaks Down}

Full fine-tuning updates \emph{all} model parameters.
For modern LLMs, this implies:
\begin{itemize}
  \item prohibitive compute cost,
  \item catastrophic forgetting,
  \item versioning and rollback complexity,
  \item high operational risk.
\end{itemize}

As a result, most production systems rely on \emph{parameter-efficient adaptation}.

\subsection{Parameter-Efficient Fine-Tuning (PEFT)}

PEFT methods freeze the base model and train a small number of additional parameters.

Benefits:
\begin{itemize}
  \item faster training,
  \item reduced memory footprint,
  \item safer updates,
  \item easier rollback.
\end{itemize}

LoRA is the dominant PEFT technique for LLMs.

\subsection{Low-Rank Adaptation (LoRA)}

LoRA modifies a linear layer:

\[
W x \rightarrow (W + \Delta W) x
\]

where:

\[
\Delta W = A B,\quad A \in \mathbb{R}^{d \times r},\ B \in \mathbb{R}^{r \times k}
\]

with rank $r \ll \min(d,k)$.

Only $A$ and $B$ are trained; $W$ remains frozen.

\subsection{Why LoRA Works}

Empirically, many adaptation tasks lie in a low-rank subspace.
LoRA exploits this by:
\begin{itemize}
  \item preserving pretrained knowledge,
  \item limiting update capacity,
  \item improving stability.
\end{itemize}

\subsection{Quantization}

Quantization reduces numerical precision (e.g.\ FP32 → FP16 → INT8).

Benefits:
\begin{itemize}
  \item reduced memory usage,
  \item faster inference,
  \item lower energy consumption.
\end{itemize}

Risks:
\begin{itemize}
  \item numerical instability,
  \item output drift,
  \item amplified rare-token errors.
\end{itemize}

\subsection{Systems Engineering Implications}

\textbf{Requirements}
\begin{itemize}
  \item Fast iteration with bounded risk
  \item Predictable inference cost
\end{itemize}

\textbf{Architecture}
\begin{itemize}
  \item Frozen base model + adapters
  \item Precision-aware serving stack
\end{itemize}

\textbf{Verification \& Validation}
\begin{itemize}
  \item Regression testing against base model
  \item Task-specific evaluation suites
\end{itemize}

\textbf{Operations}
\begin{itemize}
  \item Adapter versioning
  \item Rollback and A/B testing
\end{itemize}

\subsection{Human Factors and Failure Modes}

Common misconceptions:
\begin{itemize}
  \item ``The model has been retrained''
  \item ``Quantization only affects speed''
\end{itemize}

In reality, adaptation changes behavior subtly and requires careful validation.

\subsection{Key References}

\begin{itemize}
  \item \citet{hu2021lora} — LoRA
  \item \citet{dettmers2023qlora} — QLoRA
\end{itemize}

\subsection{Recommended University Lectures}

\begin{itemize}
  \item Stanford CS324 — model adaptation and deployment  
  \url{https://stanford-cs324.github.io/winter2022/}

  \item MIT 6.S191 — practical deep learning systems  
  \url{https://introtodeeplearning.mit.edu}
\end{itemize}

\subsection{Practitioner Lab}

See \texttt{labs/week07/week07\_lab\_adaptation\_lora\_quantization.ipynb}.

\subsection{Week 7 Takeaway}

\begin{quote}
Most real-world LLM improvements come from *adapting* models safely, not retraining them from scratch.
\end{quote}

